\renewcommand\epigraphflush{flushright}
\renewcommand\epigraphsize{\normalsize}
\setlength\epigraphwidth{0.7\textwidth}

\definecolor{titlepagecolor}{cmyk}{1,.60,0,.40}

%\DeclareFixedFont{\titlefont}{T1}{ppl}{b}{it}{0.5in}

\makeatletter                       
\def\printauthor{%                  
	{\large \@author}}              
\makeatother
\author{%
	Author 1 name \\
	Department name \\
	\texttt{email1@example.com}\vspace{20pt} \\
	Author 2 name \\
	Department name \\
	\texttt{email2@example.com}
}


\newcommand{\trtitle}{Titel Diplomarbeit}
\newcommand{\trquote}{Pure mathematics is on the whole distinctly more useful than applied. For what is useful above all is technique, and mathematical technique is taught mainly through pure mathematics.}
\newcommand{\trquoteorigin}{London 1941}
\newcommand{\trquoteman}{G. H. Hardy}
\newcommand{\trtype}{Diplomarbeit}
\newcommand{\trauthor}{Dominic \textsc{Bloech}, BSc}
\newcommand{\trauthoradress}{Abelegasse 28/5, 1160 Vienna}
\newcommand{\trstrasse}{Wiedner Hauptstraße}
\newcommand{\trmatrikelnummer}{1025929}
\newcommand{\trort}{Vienna}
\newcommand{\trbetreuer}{Ing. Dipl.-Ing.  Reinhard \textsc{Stadlmayr}, BSc.}
\newcommand{\trbetreuerberni}{Dipl.-Ing. Bernhard \textsc{Berger}, BSc.}
\newcommand{\trprof}{Univ.-Prof. Mag.rer.nat. Dipl.-Ing. Dr.techn. Friedrich \textsc{Aumayr}}
\newcommand{\trfachgebiet}{Atom- und Plasmaphysik}
\newcommand{\trinstitut}{Angewandte Physik}
\newcommand{\trfakultaet}{Arbeitsgruppe \textsc{Aumayr}}
\newcommand{\truni}{\textsc{Technischen Universitt Wien}}
\newcommand{\trunititel}{\textsc{Technische Universitt Wien}}
\newcommand{\trdate}{\today}


% The following code is borrowed from: http://tex.stackexchange.com/a/86310/10898

\newcommand\titlepagedecoration{%
	\begin{tikzpicture}[remember picture,overlay,shorten >= -10pt]
	
	\coordinate (aux1) at ([yshift=-15pt]current page.north east);
	\coordinate (aux2) at ([yshift=-410pt]current page.north east);
	\coordinate (aux3) at ([xshift=-4.5cm]current page.north east);
	\coordinate (aux4) at ([yshift=-150pt]current page.north east);
	
	\begin{scope}[titlepagecolor!40,line width=12pt,rounded corners=12pt]
	\draw
	(aux1) -- coordinate (a)
	++(225:5) --
	++(-45:5.1) coordinate (b);
	\draw[shorten <= -10pt]
	(aux3) --
	(a) --
	(aux1);
	\draw[opacity=0.6,titlepagecolor,shorten <= -10pt]
	(b) --
	++(225:2.2) --
	++(-45:2.2);
	\end{scope}
	\draw[titlepagecolor,line width=8pt,rounded corners=8pt,shorten <= -10pt]
	(aux4) --
	++(225:0.8) --
	++(-45:0.8);
	\begin{scope}[titlepagecolor!70,line width=6pt,rounded corners=8pt]
	\draw[shorten <= -10pt]
	(aux2) --
	++(225:3) coordinate[pos=0.45] (c) --
	++(-45:3.1);
	\draw
	(aux2) --
	(c) --
	++(135:2.5) --
	++(45:2.5) --
	++(-45:2.5) coordinate[pos=0.3] (d);   
	\draw 
	(d) -- +(45:1);
	\end{scope}
	\end{tikzpicture}%
}


	\begin{titlepage}
		
		\begin{flushleft}
			Unterschrift des Betreuers
		\end{flushleft}
		
		\noindent
		\textbf{\huge \trtitle}\par
		\epigraph{\trquote}%
		{\textit{\trquoteorigin}\\ \textsc{\trquoteman}}
		\null\vfill
		\vspace*{1cm}
		\noindent
		\hfill
		
		
		\begin{center}
			\textbf{\trtype} \\
			im Fachgebiet \trfachgebiet \\
			vorgelegt bei \trprof \\
			am Institut für \trinstitut \\
			in der \trfakultaet \\
			an der \truni \\[0.5cm]
			vorgelegt von \\
			\textbf{\trauthor}
			
%			\hspace{0.3cm}
			
			\begin{center}
				\begin{tabular}{ll}
					Betreut durch:& \trbetreuer  \\
					& \trbetreuerberni \\
				\end{tabular}
			\end{center}
		\end{center}
		
		\begin{flushright}
		\begin{minipage}{0.35\linewidth}
			\begin{flushright}
				\begin{tabular}{l}
					\trauthor \\
					\trauthoradress \\
					Matrikelnummer:  \trmatrikelnummer \\
					\trort, \trdate
				\end{tabular}
			\end{flushright}
		\end{minipage}
		%
		\begin{minipage}{0.02\linewidth}
			\rule{1pt}{70pt}
		\end{minipage}
		\end{flushright}
	
	
	\begin{tabular}{lcr}
		\includegraphics[scale=0.15]{TU_Logo.pdf} & % dein_unilogo.jpg/.eps im Verzeichnis "bilder" ablegen
		\hspace{1cm} \trunititel \hspace{2.5cm} &
		%\includegraphics[scale=0.3]{iap-logo.jpg} % dein_fglogo.jpg/.eps im Verzeichnis "bilder" ablegen, Fachgebietslogo
		\\
	\end{tabular}

\begin{flushleft}
	Unterschrift des Studenten
\end{flushleft}
	
		\titlepagedecoration
	\end{titlepage}

